% ===== tr-01 块5：第 6 章 pp.5--8（原文页 5--8）=====
% [原文 p.5]
平均曲率流（MCF）产生一族随时间变化的浸入 $F_t:M\subseteq X$，满足方程
\[
\frac{\partial}{\partial t}F_t=\vec H.
\]
Smoczyk 证明了在 Kähler--Einstein 流形中 MCF 保持拉格朗日条件 [58]，因此这种情形下它被称为拉格朗日平均曲率流（LMCF）。若能证明该流光滑收敛，则 LMCF 的极限满足 $\nabla\theta=0$，从而是特殊拉格朗日。然而 Neves 的例子表明，对特殊拉格朗日作任意小的 $C^0$ 形变后，其 LMCF 可能产生有限时间奇点 [56]。更糟的是，在 Tian--Yau 空间的设定下，由于单射半径退化，几何不是有界的，因此文献中的许多标准工具都不适用。尽管如此，作者得到了如下定理：

\begin{theorem}{6.2.1}
（见 [10, 定理 1.1]）对 $\varepsilon\ll 1$，$M_\varepsilon$ 的 LMCF 光滑收敛到一个光滑的特殊拉格朗日子流形。
\end{theorem}
% [Tutor Note: MCF 方程 $\frac{\partial}{\partial t}F_t=\vec H$ 中的 $\vec H$ 是平均曲率向量（沿法向），不是标量平均曲率；LMCF 的不动点满足 $\vec H=0$，即特殊拉格朗日（相为常数）。]

特别地，这在完备 Calabi--Yau 空间中产生许多紧特殊拉格朗日子流形的新例子。

\subsection{del Pezzo 曲面中的特殊拉格朗日纤维化}
本节我们限于 $Y$ 是弱 del Pezzo 曲面的情形，并解释 $X$ 中特殊拉格朗日纤维化的存在性。由紧复曲面的分类，弱 del Pezzo 曲面或者是 $\PP^1\times\PP^1$、Hirzebruch 曲面 $\mathbb{F}_2$，或者是 $\PP^2$ 在 $d\le 8$ 个点处的逐次爆破（且没有点落在 $(-2)$-曲线上）。

\subsubsection{无穷远附近的局部纤维化}
当 $Y$ 是弱 del Pezzo 曲面时，光滑除子 $D\in|-K_Y|$ 必然是椭圆曲线，从而对任意本原类 $\gamma\in H_1(D,\Z)$，它容许以 $\gamma$ 为纤维类的特殊拉格朗日纤维化。由上一节的讨论，这样的特殊拉格朗日纤维化提升为 $X_{\C}$ 中的特殊拉格朗日纤维化，并在 $X$ 的无穷远处附近产生一个几乎特殊拉格朗日的纤维化。此外，由定理 6.2.1，LMCF 把每个这样的纤维流到特殊拉格朗日环面。先验地，并不清楚 LMCF 是否保持纤维化结构------例如，考虑哑铃形时容易看出，即使是曲线缩短流也可能把互不相交的、几乎极小的子流形演化到同一个极小极限。然而在我们的情形可以利用辛结构：取 $X$ 中两个互不相交的几乎特殊拉格朗日 $M_\varepsilon,M'_{\varepsilon'}$，并假设它们在 LMCF 下收敛到特殊拉格朗日 $\tilde M_\varepsilon=\tilde M'_{\varepsilon'}$。由于 LMCF 保持拉格朗日子流形的 Hamiltonian 同痕类，而拉格朗日之间的 Floer 上同调不依赖于 Hamiltonian 同痕类中代表元的选取，我们有
\[
0=HF(M_\varepsilon,M'_{\varepsilon'})\cong HF(\tilde M_\varepsilon,\tilde M_{\varepsilon'})\cong HF(\tilde M_\varepsilon)\cong H^*(T^2)\neq 0,
\]
% [原文 p.6]
这是一个矛盾。\footnote{这里的背景几何不是有界几何，因此要定义 Floer 上同调，需要此设定下 Sachs--Uhlenbeck--Gromov 紧性定理的一个版本，它由 [10, 命题 5.3] 或 [25] 提供。}
% [Tutor Note: 最后一步 $HF(\tilde M_\varepsilon)\cong H^*(T^2)$ 是因为两个假定收敛的特殊拉格朗日 $\tilde M_\varepsilon,\tilde M'_{\varepsilon'}$ 在拓扑上相交数为零，而 $HF$ 与通常上同调一致；$H^*(T^2)\neq 0$ 与 $0=HF$ 矛盾，从而 $\tilde M_\varepsilon\neq\tilde M'_{\varepsilon'}$。]

这里我们用到了如下事实：在超 Kähler 曲面中，由维数原因，Maslov 指标为零的拉格朗日关于一般的殆复结构不界定任何 $J$-全纯圆盘，因此 Floer 上同调退化为通常的上同调。这表明 $\tilde M_\varepsilon\neq\tilde M'_{\varepsilon'}$。另一方面，经过适当的超 Kähler 旋转（如下所述），$\tilde M_\varepsilon,\tilde M'_{\varepsilon'}$ 成为全纯的。特别地，它们只能有有限多个交点，且每个交点都贡献正值。由于它们的拓扑相交数为零，这实际上表明 $\tilde M_\varepsilon,\tilde M'_{\varepsilon'}$ 互不相交。因此，LMCF 把几乎特殊拉格朗日纤维化流成无穷远附近真正的特殊拉格朗日纤维化。

\subsubsection{超 Kähler 旋转与紧化}
下面我们解释超 Kähler 旋转技巧。由于 $Sp(1)\cong SU(2)$，每个 Calabi--Yau 曲面都是超 Kähler 的。给定带全纯体积形式 $\Omega$ 与 Ricci 平坦度量 $\omega$、且满足 $2\omega^2=\Omega\wedge\bar\Omega$ 的 Calabi--Yau 曲面 $X$，可以对同一个底流形 $X$ 联系一族 $S^2$ 参数化的复结构以及相应的 Ricci 平坦 Kähler 形式。为简单起见，我们取定 $S^2$ 中一个带坐标 $\theta$ 的 $S^1$，其相应的 Kähler 形式与全纯体积形式由
\[
\omega_\theta=\operatorname{Re}(e^{-i\theta}\Omega),\qquad \Omega_\theta=\omega-i\operatorname{Im}(e^{-i\theta}\Omega)
\]
给出。我们记 $X_\theta$ 为配备 Kähler 形式 $\omega_\theta$ 与全纯体积形式 $\Omega_\theta$ 的空间 $X$。由 Wirtinger 定理，$C\subseteq X$ 是全纯曲线，当且仅当 $C$ 在同一个底空间上、但带有由 $X_\theta$ 给出的 Calabi--Yau 结构时是特殊拉格朗日的。特别地，对超 Kähler 曲面，全纯闭链与特殊拉格朗日闭链之间有密切的联系。
% [Tutor Note: 这里 $X_\theta$ 与 $X$ 是同一个 Riemann 流形（同一个底空间），只是复结构不同；$\omega_\theta,\Omega_\theta$ 由 $\Omega$ 的相位旋转给出，这就是"超 Kähler 旋转"的含义------不能理解为不同的流形。]
% [Tutor Note: $2\omega^2=\Omega\wedge\bar\Omega$ 是曲面上 Calabi--Yau 结构的归一化条件（$\Omega$ 为 $(2,0)$-形式），它保证 $\omega_\theta$ 仍是 Kähler 形式；$\operatorname{Re}$、$\operatorname{Im}$ 为正体算子记号。]

现在，由定理 6.2.1 与上述讨论，我们在 $X$ 的无穷远处附近构造了一个特殊拉格朗日环面纤维化；经过适当的超 Kähler 旋转后，它变成一个亏格一的全体纯纤维化。我们把这个超 Kähler 旋转记作 $\check X$，其 Kähler 形式 $\check\omega$ 与全纯体积形式 $\check\Omega$ 由
\begin{equation}
\check\omega=\operatorname{Re}\Omega,\qquad \check\Omega=\omega-i\operatorname{Im}\Omega
\tag{6.2.2}
\end{equation}
给出。
% [Tutor Note: 注意式 (6.2.2) 中两个式子的角色：$\check\omega$ 取 $\Omega$ 的实部，而 $\check\Omega$ 由 $\omega$ 与 $\operatorname{Im}\Omega$ 组合而成。不要误写成 $\check\omega=\omega$、$\check\Omega=\Omega$。]

不难看出纤维化存在局部截面。因此存在紧集 $K\subseteq\check X$，使得 $\check X\setminus K\to\Delta^*$ 是椭圆纤维化，其中 $\Delta^*$ 是去心圆盘。特别地，存在从 $\Delta^*$ 到带标记点的椭圆曲线模空间 $\mathcal{M}_{1,1}$ 的全纯映射。由于纤维化绕无穷远的拓扑单值与 $\begin{pmatrix}1&d\\0&1\end{pmatrix}$ 共轭，椭圆纤维化 $\check X\setminus K\to\Delta^*$ 是 $\mathcal{M}_{1,1}$ 上泛族经由 [29, 命题 5.9] 的拉回。此外，
% [原文 p.7]
$\check X$ 中无穷远处附近椭圆纤维化的 $j$-不变量趋于无穷，而全纯映射 $\Delta^*\to\mathcal{M}_{1,1}$ 延拓为 $\Delta\to\bar{\mathcal{M}}_{1,1}$。把 $\bar{\mathcal{M}}_{1,1}$ 上的泛族拉回，得到 $\check X$ 的一个紧化 $\check Y$。由纤维化绕无穷远的单值，在无穷远处添加的纤维 $\check D$ 属于 Kodaira 分类中的 $I_d$ 型。总之，$X$ 的一个合适的超 Kähler 旋转 $\check X$ 可以通过添加一条 $I_d$-纤维而紧化成一个紧复曲面。
% [Tutor Note: 单值矩阵 $\begin{pmatrix}1&d\\0&1\end{pmatrix}$ 中的 $d$ 就是 Kodaira 纤维型 $I_d$ 的 $d$；这与 Intro 中 del Pezzo 曲面的次数 $d$ 是同一个 $d$（因为 $\check Y$ 的紧化恰好添加 $I_d$ 纤维）。]
% [Tutor Note: $\mathcal{M}_{1,1}$ 是带一个标记点的椭圆曲线模空间（开），$\bar{\mathcal{M}}_{1,1}$ 是它的紧化（Deligne--Mumford 紧化，带上尖点对应 $I_d$ 型退化）；原文两处都带/不带横线，需区分。]

\subsection{Enriques--Kodaira 分类的推论}
现在我们可以利用紧复曲面的 Enriques--Kodaira 分类来推出 $\check Y$ 的结构。回忆如下定义 [31]：

\begin{definition}{6.2.2}
有理椭圆曲面 $\check Y$ 是一个光滑紧复曲面，满足以下任一等价定义：
\begin{enumerate}
\item[(i)] $\check Y$ 是带有一个截面的椭圆纤维化结构的有理曲面。
\item[(ii)] $\check Y$ 可以实现为 $\PP^2$ 中三次曲线束的基轨迹的爆破。
\item[(iii)] $-K_{\check Y}$ 无基点，并定义一个纤维化。
\end{enumerate}
\end{definition}

设 $\check\omega,\check\Omega$ 是如上由超 Kähler 旋转得到的 $\check X$ 上相应的 Ricci 平坦 Kähler 形式与全纯体积形式。由局部计算可以验证，$\check\Omega$ 延拓为 $\check Y$ 上的一个亚纯 2-形式，它沿 $\check D$ 有单极点，其中 $\check D=\check Y\setminus\check X$。因此我们有 $K_{\check Y}=\mathcal{O}_{\check Y}(-\check D)$。由 $\check Y$ 上的椭圆纤维化可得 $c_1(\check Y)^2=0$。由伴随公式，纤维中没有 $(-1)$-曲线。由于 del Pezzo 曲面或者是 $\PP^2$ 的爆破、或者是 $\PP^1\times\PP^1$，我们有 $b_1(\check Y)=0$。由 Mayer--Vietoris 序列，我们有
\[
b_1(\check Y)=b_1(\check X)=b_1(X)=b_1(Y)=0.
\]
% [Tutor Note: 这一串等式来自 $X=Y\setminus D$ 与 $\check X=\check Y\setminus\check D$ 的配对（Mayer--Vietoris），以及 $X$ 与 $\check X$ 通过超 Kähler 旋转是同胚的；因此四个第一 Betti 数相等。]
% [Tutor Note: 定义 6.2.2 的 (ii) 中"三次曲线束的基轨迹"（base locus of a pencil of cubics）是 $\PP^2$ 中 9 个交点的轨迹，其爆破给出有理椭圆曲面；这与 (i)(iii) 的等价性是有理椭圆曲面理论的标准结论。]

若 $\check Y$ 极小，则由 Enriques--Kodaira 分类（例如见 [3, 第 VI 章, 表 10]），$\check Y$ 只能是 Enriques 曲面、K3 曲面或极小真椭圆曲面。由于真椭圆曲面的 Kodaira 维数为 1，而 $K_{\check Y}=\mathcal{O}_{\check Y}(-\check D)$ 这一事实排除了全部三种可能，因此 $\check Y$ 不是极小的。

% [原文 p.8]
现在，$\check Y$ 中任一 $(-1)$-曲线 $E$ 与 $\check D$ 的相交数为 1，因此 $(\check D+E)^2>0$。所以由 [3, 第 IV 章, 定理 5.2]，$\check Y$ 是射影的。于是由 Hodge 理论有 $h^1(\check Y,\mathcal{O}_{\check Y})=0$，且由于 $-K_{\check Y}$ 有效，有 $h^0(\check Y,K_{\check Y}^2)=0$。最后，Castelnuovo 的有理性判别法蕴含 $\check Y$ 是有理的。再由 Riemann--Roch 定理得到 $H^0(\check Y,\mathcal{O}_{\check Y}(\check D))=2$，而 $\check Y$ 中 $\check D$ 附近的局部椭圆纤维化实际上延拓为一个椭圆纤维化。确实，由于 $H^1(\check Y,\mathcal{O}_{\check Y})=H^2(\check Y,\mathcal{O}_{\check Y})=0$，有 $\operatorname{Pic}(\check Y)\cong H^2(\check Y,\Z)$。因此 $\check Y$ 是有理椭圆曲面。特别地，存在椭圆纤维化 $\check X\to\C$。反转超 Kähler 旋转，就得到特殊拉格朗日纤维化 $X\to\R^2$。或者，也可以直接由形变与 $J$-全纯曲线的紧性理论论证 $X\setminus K$ 上的纤维化延拓为整个 $X$ 上的全纯纤维化；见 [10]。总结起来，我们有如下定理
% [Tutor Note: $( \check D+E )^2>0$ 中 $(-1)$-曲线 $E$ 满足 $E^2=-1$、$E\cdot\check D=1$，展开得 $\check D^2+2\check D\cdot E+E^2>0$；这是证明 $\check Y$ 射影性的关键一步（Nakai--Moishezon 判别法）。]

\begin{theorem}{6.2.3}
（见 [10, 定理 1.3、定理 1.6]）设 $Y$ 是次数为 $d$ 的 del Pezzo 曲面，$D$ 是光滑反典范除子，则 $X=Y\setminus D$ 容许特殊拉格朗日纤维化。此外，由式 (6.2.2) 给出的 $X$ 的合适超 Kähler 旋转 $\check X$ 可以通过添加一条 $I_d$ 纤维而紧化为有理椭圆曲面 $\check Y$。
\end{theorem}

\begin{remark}{6.2.4}
在 $Y=\PP^2$ 的情形，这证明了 Auroux 的一个猜想 [1, 猜想 5.1]。值得指出的是，相应的有理椭圆曲面 $\check Y$ 具有奇异构型 $I_9I_1^3$。对镜像对称感兴趣的读者不难把它与 $\PP^2$ 的镜像超势联系起来。$\PP^2$ 的镜像超势是 $x+y+\frac{1}{xy}:(\C^*)^2\to\C$，而它恰好紧化为具有奇异构型 $I_9I_1^3$ 的唯一极端有理椭圆曲面。这暗示了镜像对称与超 Kähler 旋转之间的联系；进一步讨论见第 6.6.1 节。
\end{remark}
% [Tutor Note: $I_9I_1^3$ 表示一条 $I_9$ 型奇异纤维与三条 $I_1$ 型奇异纤维（共 12 条奇异纤维，与椭圆曲面的 Euler 数 12 相符）；不要读成 $I_9 I_1^3$ 中的指数混淆------$I_1$ 的立方即三条 $I_1$。]

\begin{remark}{6.2.5}
当 $Y$ 是有理椭圆曲面、$D$ 是光滑椭圆纤维时，Tian--Yau [62] 构造了一个渐近柱状的 Ricci 平坦度量 [62, 定理 5.2]。定理 6.2.3 中的论证在此情形也适用，见 [10, 定理 1.3]。特别地，对一般的 $(Y,D)$ 选取，$X$ 中的特殊拉格朗日纤维化有 12 条 $I_1$ 型奇异纤维。这证实了 Auroux 的另一个猜想 [1, 猜想 2.10]。
\end{remark}

\begin{remark}{6.2.6}
回忆：每个本原类 $\gamma\in H_1(D,\Z)$ 生成 $X$ 中的一个特殊拉格朗日纤维化，其纤维的同调类等于 $S^1$ 与 $D$ 中同调类为 $\gamma$ 的特殊拉格朗日之积。特别地，$X$ 中有可数多个不同的特殊拉格朗日纤维化结构。
\end{remark}

\section{有理椭圆曲面上的特殊拉格朗日纤维化}
本节我们讨论有理椭圆曲面上的 Calabi--Yau 度量与特殊拉格朗日纤维化。设 $\check Y\to\PP^1$ 是如定义 6.2.2 所述的有理椭圆曲面。由爆破的典范丛公式，椭圆纤维化的每条纤维都是反典范除子。取任一纤维 $\check D\subseteq\check Y$，并令 $\check X=\check Y\setminus\check D$，则 $\check Y$ 上存在一个沿 $\check D$ 有单极点的亚纯形式 $\check\Omega$，它在 $\C^*$-伸缩下唯一。考虑对 $(\check Y,\check X)$ 的整系数长正合序列，我们有
\begin{equation}
0\to H^1(\check D)\to H_2(\check X)\to H_2(\check Y)\to H^2(\check D).
\tag{6.3.1}
\end{equation}
其中我们用到了 Poincaré 对偶 $H^k(\check D)\cong H_{4-k}(\check Y,\check X)$。设 $[C]\in H_2(\check X)$ 是 $H^1(\check D)\cong\Z$ 的生成元的像，它在 $\check X$ 中被称为"坏轮"（bad cycle）。
% [Tutor Note: (3.1)（本译本式 (6.3.1)）是同调/上同调混合的长正合序列：$H^1(\check D)$ 与 $H^2(\check D)$ 是上同调，中间的 $H_2$ 是同调；用 Poincaré 对偶 $H^k(\check D)\cong H_{4-k}(\check Y,\check X)$ 把配对 $(\check Y,\check X)$ 的同调统一起来。]
% [Tutor Note: "坏轮"（bad cycle）$[C]$ 是 $H^1(\check D)$ 生成元的像；它是唯一一个推到 $H_2(\check Y)$ 中同调于零的类，这正是它"坏"（阻碍 Ricci 平坦度量的存在）的几何来源。]
